\section{Related Work}
\label{sec:related}
% Every \cite key must exist in refs.bib AND pass the project's bibliography verification (scripts/verify_bib.py).
% Do not add a citation in Overleaf without sending it back for verification first: citing non-existent work is a desk-reject reason.
% \todo for co-authors: candidates still to verify before citing: SeekVLN's companion benchmarks (e.g., ESI-Bench, EXPRESS-Bench).

\paragraph{RL for VLM agents and evidence seeking.}
VAGEN~\cite{vagen2025} trains VLM agents with multi-turn RL and world-model reasoning rewards; we build on its trainer.
GLANCE~\cite{glance2026} turns the agent's linguistic predictions into a curiosity signal by comparing predicted and observed visual representations.
Closest to us, concurrent work on SeekVLN~\cite{seekvln2026} observes that VLM navigation agents act confidently when key evidence is missing, and branches from the same state into evidence seeking and direct navigation to credit the seeking decision by the subsequent return.
Its credit is computed in the realized world and the world itself never changes; our pairs average credit over the worlds consistent with the agent's information and additionally supervise whether the goal should change after looking.
We include a single-world look-versus-commit contrast in the spirit of that work as one cell of our $2{\times}2$ design.
Asking for help~\cite{askforhelp2025} and compare-before-commit navigation~\cite{c2nav2026} address related decisions about when to act on uncertain information.

\paragraph{Counterfactual data and supervision.}
Counterfactual VLN~\cite{counterfactualvln2020} trains navigation agents to behave consistently across original and counterfactual environments.
Grounding on Request~\cite{groundingonrequest2026} uses paired interventions and a twin-margin objective to teach how answers should change when the requested evidence changes and when they should not; its paired training uses given inputs, whereas our consistency term acts only after the agent has acquired the evidence through its own actions.
DEER-3D~\cite{deer3d2025} edits 3D scenes to produce counterfactual question--answer pairs for grounding, and DUEL~\cite{duel2026} perturbs textual statements in a self-play loop; both are static and do not involve acquiring evidence.
Several works study how much static VLM reasoning depends on the image~\cite{thinkingdeltas2026,morethinking2025}, and language interventions show that vision-language-action models often execute the original goal~\cite{encodednotcontrol2026}.

\paragraph{Environment design and robustness.}
Unsupervised environment design~\cite{paired2020,accel2022}, LLM-driven environment designers~\cite{spade2026,simworldstudio2026}, quality-diversity search for failures of VLA policies~\cite{qdeg2026}, physical red-teaming~\cite{redvla2026}, and perturbation benchmarks~\cite{liberoplus2026} generate or search for environments that challenge an agent.
We do not search for hard worlds; we construct pairs whose relation to the agent's information is known and verified, and use that relation in the learning signal.

\paragraph{Partial observability and variance reduction.}
Learning counterfactual world models under partial observability~\cite{clwm2026} studies when action sequences have distinguishable consequences.
Our pre-evidence credit is a stratified estimator over the hidden world combined with a commit-now baseline; it is related to classical variance reduction in policy gradients~\cite{gae2016}, but applies it to evidence-seeking decisions with verified counterfactual worlds.
